\documentclass[12pt]{article}

% Packages
\usepackage[utf8]{inputenc}
\usepackage[T1]{fontenc}
\usepackage{lmodern}   % Use Latin Modern (Computer Modern-like) font
\usepackage{amsmath, amssymb, amsthm}
\usepackage{graphicx}
\usepackage{hyperref}
\usepackage{geometry}
\usepackage{caption}
\usepackage{subcaption}
\usepackage[numbers]{natbib}
\usepackage{setspace}
\usepackage{booktabs}
\usepackage{float}
\usepackage{bm}
\usepackage{listings}
\usepackage{xcolor}

% Page Setup
\geometry{a4paper, margin=1in}
\onehalfspacing

% Code style
\lstdefinestyle{pycode}{
  language=Python,
  basicstyle=\ttfamily\small,
  keywordstyle=\bfseries\color{blue},
  commentstyle=\itshape\color{gray!70!black},
  showstringspaces=false,
  columns=fullflexible,
  frame=single,
  framerule=0.5pt,
  breaklines=true
}

% Title Information
\title{\textbf{Gradient-Based Sensitivity Analysis in Large Language Models}}
\author{}
\date{}

% Begin Document
\begin{document}

\maketitle
\vspace{-1.5em}
\hrulefill

\begin{abstract}
We formalize and implement a gradient-based sensitivity analysis for neural language models. 
Theoretical connections between per-parameter gradients and sensitivity to soft errors motivate a practical procedure: 
stream a language modeling corpus into fixed windows, compute next-token cross-entropy loss, backpropagate, 
and track the maximum absolute gradient observed per weight element across many batches to identify the most sensitive parameters.
\end{abstract}

\section{Gradient-Based Sensitivity Framework}
Certain parameters (elements of weight tensors) are more impactful than others when a soft error is introduced. 

\paragraph{Notation.} We define the following:
\begin{itemize}
  \item $\vec{\theta} \in \mathbb{R}^N$: The parameter vector
  \[
    \vec{\theta} =
    \begin{pmatrix}
    \theta_{1} \\
    \theta_{2} \\
    \vdots \\
    \theta_{N}
    \end{pmatrix}.
  \]
  \item $x$: The model input (e.g., a token sequence).
  \item $y$: The target/label for the input (e.g., the true next token).
  \item $f(x;\vec{\theta})$: The LLM's output when parameters are $\vec{\theta}$.
  \item $\mathcal{L}(f(x;\vec{\theta}),y)$: The scalar loss function (e.g., cross-entropy) measuring mismatch between prediction and target.
  \item $\delta$: A small scalar perturbation magnitude.
  \item $e_i$: The $i$-th standard basis vector with a $1$ in the $i$-th position and $0$ elsewhere.
\end{itemize}

\paragraph{Sensitivity.}
The sensitivity of parameter $\theta_i$ is defined as
\begin{equation}
  S_i(\delta) \;=\; \mathcal{L}\!\big(f(x;\vec{\theta}+\delta e_i),y\big) \,-\, \mathcal{L}\!\big(f(x;\vec{\theta}),y\big).
\end{equation}
Using a first-order Taylor expansion, this can be approximated as
\begin{equation}
  S_i(\delta) \;\approx\; \delta \, g_i,
  \qquad \text{where}\quad \bm{g} = \nabla_{\vec{\theta}}\,\mathcal{L}(f(x;\vec{\theta}),y).
\end{equation}
Thus, the gradient entry $g_i$ serves as a first-order proxy for sensitivity.

\section{Methodology and Procedure}

\paragraph{Corpus as a vector.} 
The tokenized corpus is represented as
\[
x = [x_0, x_1, \ldots, x_{T-1}], \qquad x_t \in \{0,1,\ldots,V-1\},
\]
where $T$ is the total number of tokens and $V$ is the vocabulary size.

\paragraph{Windowing.} 
For each index $n$, we extract a window
\[
w(n) = [x_n, x_{n+1}, \ldots, x_{n+L}],
\]
where $L = \texttt{SEQ\_LEN}$ is the input sequence length.

\paragraph{Inputs and labels.} 
Each window is split into input and label sequences:
\[
I(n) = [x_n, \ldots, x_{n+L-1}], \qquad 
Y(n) = [x_{n+1}, \ldots, x_{n+L}].
\]

\paragraph{Batching.} 
Stacking $B$ such pairs yields input and label tensors
\[
I \in \mathbb{N}^{B \times L}, 
\qquad 
Y \in \mathbb{N}^{B \times L}.
\]

\paragraph{Forward pass.} 
The model with parameters $\vec{\theta}$ produces logits
\[
Z = f(I;\vec{\theta}) \in \mathbb{R}^{B \times L \times V}.
\]

\paragraph{Loss.} 
The cross-entropy loss averaged over batch and sequence positions is
\[
\mathcal{L}(I,Y;\vec{\theta}) 
= \frac{1}{B(L-1)} \sum_{b=1}^B \sum_{t=0}^{L-2} 
    -\log p\!\left(Y_{b,t+1}\,\middle|\,I_{b,\leq t}; \vec{\theta}\right).
\]

\paragraph{Gradient.} 
By backpropagation, the gradient of the loss with respect to the parameters is
\[
\nabla_{\vec{\theta}} \, \mathcal{L}(I,Y;\vec{\theta})
= \begin{bmatrix}
\frac{\partial \mathcal{L}}{\partial \theta_1},\;
\frac{\partial \mathcal{L}}{\partial \theta_2},\;
\ldots,\;
\frac{\partial \mathcal{L}}{\partial \theta_N}
\end{bmatrix}^{\!\top}.
\]

This gradient vector assigns one sensitivity score per parameter. 
The magnitude 
\[
\left| \frac{\partial \mathcal{L}}{\partial \theta_i} \right|
\]
highlights the most sensitive elements in the model.

\section{Experimental Results}
The above procedure was conducted on the smallest model of GPT-2 (137m parameters). A sequence length of 1024 tokens, batch size of 16 and step size of 1100 was used to cover $\sim$18 million tokens (approx. 18\%) of the Wikitext-103 text corpus. 

\subsection*{Top-10 Most Sensitive Parameters}
The top-10 most sensitive parameters were found to be:

\begin{table}[H]
\centering
\begin{tabular}{|c|c|c|c|}
\hline
Rank & Weight & Element & $|\nabla \mathcal{L}|$ \\
\hline
1 & transformer.wte.weight & (2488,496)  & 5.235e+00 \\
2 & transformer.wte.weight & (837,496)   & 4.481e+00 \\
3 & transformer.wte.weight & (198,496)   & 3.247e+00 \\
4 & transformer.wte.weight & (11,496)    & 2.898e+00 \\
5 & transformer.wte.weight & (34315,496) & 2.669e+00 \\
6 & transformer.wte.weight & (796,496)   & 2.517e+00 \\
7 & transformer.wte.weight & (764,496)   & 2.499e+00 \\
8 & transformer.wte.weight & (13,496)    & 2.427e+00 \\
9 & transformer.wte.weight & (220,496)   & 2.242e+00 \\
10 & transformer.wte.weight & (2488,430) & 2.207e+00 \\
\hline
\end{tabular}
\end{table}

The Top-10 Most Sensitive Parameters (excluding the embedding layer, and ayer norm weights):

\begin{table}[H]
\centering
\begin{tabular}{|c|c|c|c|c|}
\hline
Local Rank & Global Rank & Weight & Element & $|\nabla \mathcal{L}|$ \\
\hline
1 & 117 & transformer.h.5.mlp.c\_fc.weight & (393, 1866) & 1.341e+00 \\
2 & 123 & transformer.h.5.mlp.c\_fc.weight & (373, 1866) & 1.319e+00 \\
3 & 134 & transformer.h.5.mlp.c\_fc.weight & (64, 1866)  & 1.298e+00 \\
4 & 158 & transformer.h.5.mlp.c\_fc.weight & (447, 1866) & 1.177e+00 \\
5 & 219 & transformer.h.5.mlp.c\_fc.weight & (326, 1866) & 1.007e+00 \\
6 & 242 & transformer.h.5.mlp.c\_fc.bias   & (1866)      & 9.751e-01 \\
7 & 243 & transformer.h.5.mlp.c\_fc.weight & (138, 1866) & 9.741e-01 \\
8 & 246 & transformer.h.6.attn.c\_attn.weight & (64, 2197)  & 9.673e-01 \\
9 & 262 & transformer.h.7.attn.c\_attn.weight & (64, 2296)  & 9.316e-01 \\
10 & 300 & transformer.h.3.mlp.c\_fc.weight & (393, 1612) & 8.853e-01 \\
\hline
\end{tabular}
\end{table}

\subsection*{Bit-Flip Trials}
A single bit flip was conducted repeatedly (e.g. 5 times) on different sections of the bits (sign, exponent, mantissa bits) for different input prompts. The following prompts were taken:

\begin{enumerate}
\item The weather today is
\item The patient should take
\item The bank transfer amount is
\item The recommended dose for a child is
\item The evacuation order status is
\end{enumerate}

The model generated 100 tokens as the output. The outputs were classified into three categories:

\begin{itemize}
\item \textbf{Preserved} - The meaning of the output was relatively unchanged. Condition: BERTScore\_F1 $>$ 0.87
\item \textbf{Changed} - The output was intelligible, but the meaning had changed. Condition: 0.80 $<$ BERTScore\_F1 $<$ 0.87
\item \textbf{Gibberish} - The output had completely become a random jumble of words and symbols. Condition:  BERTScore\_F1 $<$ 0.80
\end{itemize}

\subsection*{Top-3 Most Sensitive Element Test}
The top-3 most sensitive tensor elements were tested. Out of 225 outputs:

\begin{table}[H]
\centering
\begin{tabular}{|c|c|p{8cm}|}
\hline
Classification & Count & Example \\
\hline
Preserved & 4 & Input: The patient should take \newline
Clean Output: a blood test to confirm the diagnosis. The patient should be given a blood test to confirm the diagnosis. \newline
Corrupted Output: a blood test to confirm the diagnosis. The patient will be given a blood sample and a urine sample. \\
\hline
Changed & 6 & Input: The weather today is \newline
Clean Output: very good, and we're going to be able to get some good weather tomorrow. \newline
Corrupted Output: very good, and we're going to be able to get some good weather tomorrow. "We're going into the weekend and we'll be able, hopefully, to get a good start to the season." The Reds have been without a win in their last five games, with a 3-1 defeat at home to the Royals on Saturday..@RedsReds are in the midst of a 3rd straight winless start. \#Reds — RedsNation (@R \\
\hline
Gibberish & 215 & Input: The evacuation order status is \newline
Clean Output: being reviewed by the Department of Homeland Security. The evacuation order was issued by the Department of Homeland Security on July 1, 2017. \newline
Corrupted Output: , , , and , , respectively , , . , ,, , or , , are , , ( , , ), , , etc , , in , , as , , \& , , for , , with , , without , , by , , on , , under , , out , , to , , through , , over , , at , , from , , into , , back , , across , , between , , within , , around , , outside , , inside , , above \\
\hline
\end{tabular}
\end{table}

\subsection*{Analysis}
The results indicate that the vast majority ($95.6\%$) of outputs degraded into gibberish when a single bit flip was introduced in the top-3 most sensitive parameters, which were classified as 'Gibberish'. A smaller fraction ($2.67\%$) produced intelligible text whose semantic meaning diverged significantly from the clean baseline, which were classified as 'Changed', while only $1.78\%$ of generations maintained their intended meaning, which were classified as 'Preserved'. These results underscore both the extreme fragility of highly sensitive parameters to low-level perturbations and the rarity of semantically robust outputs under such faults.

\section{Conclusion}
We presented a gradient-based framework for estimating parameter sensitivity and a practical scan that identifies the most sensitive elements in large language models. 
This offers a principled way to target robustness interventions (e.g., bit-flip testing, weight clipping, or redundancy) at the most impactful parameters.

\appendix

\section{Tokenization and Batching Procedure}
The dataset is tokenized and streamed into a rolling buffer (cache). From this buffer, windows of fixed length are extracted, split into inputs and labels, and grouped into batches.

\subsection*{Algorithm A.1: Tokenization and Batching}
\textbf{Input:} Tokenized corpus $\mathbf{x} = [x_0, x_1, \ldots, x_{T-1}]$, sequence length $L$, batch size $B$. \\
\textbf{Output:} Batches of inputs $I \in \mathbb{N}^{B \times L}$ and labels $Y \in \mathbb{N}^{B \times L}$.

\begin{enumerate}
  \item Initialize an empty rolling buffer \texttt{cache}.
  \item Stream tokens $x_t$ into the buffer.
  \item While the buffer has at least $L+1$ tokens:
  \begin{enumerate}
    \item Take a window $\mathbf{w} = [x_n, \ldots, x_{n+L}]$.
    \item Define inputs $\mathbf{I} = [x_n, \ldots, x_{n+L-1}]$.
    \item Define labels $\mathbf{Y} = [x_{n+1}, \ldots, x_{n+L}]$.
    \item Store $(\mathbf{I}, \mathbf{Y})$.
    \item Remove consumed tokens from the buffer.
  \end{enumerate}
  \item Group stored pairs into mini-batches of size $B$.
\end{enumerate}

\subsection*{Code Illustration}
\begin{lstlisting}[style=pycode]
def chunk_generator(x, L):
    cache = []
    for token in x:               # stream tokens
        cache.append(token)
        while len(cache) >= L + 1:
            win, cache = cache[:L+1], cache[L+1:]
            inp, labels = win[:-1], win[1:]
            yield inp, labels

def get_batches(x, L, B):
    gen = chunk_generator(x, L)
    buf = []
    for inp, labels in gen:
        buf.append((inp, labels))
        if len(buf) == B:
            I = torch.tensor([p[0] for p in buf], dtype=torch.long)
            Y = torch.tensor([p[1] for p in buf], dtype=torch.long)
            yield I, Y
            buf = []
\end{lstlisting}

\section{Gradient Scan Procedure}
After the batches are constructed, they are passed through the model to compute the loss. PyTorch’s autograd engine then computes gradients for all parameters during backpropagation. To identify the most sensitive parameters, we record the maximum absolute gradient per element over many batches.

\subsection*{Algorithm B.1: Gradient Scan}
\textbf{Input:} Model $f(\cdot;\vec{\theta})$, batches $(I,Y)$, maximum steps $S$. \\
\textbf{Output:} For each parameter element $\theta_i$, the maximum magnitude of the observed gradientradient.

\begin{enumerate}
  \item Initialize a dictionary \texttt{running\_max} mapping each parameter tensor to zeros of its shape.
  \item For each batch $(I,Y)$ up to $S$ steps:
  \begin{enumerate}
    \item Zero model gradients.
    \item Compute the forward pass: $Z = f(I;\vec{\theta})$.
    \item Compute the loss: $\mathcal{L}(I,Y;\vec{\theta})$.
    \item Backpropagate: \texttt{loss.backward()}.
    \item For each parameter tensor $p$:
    \begin{equation*}
      \texttt{running\_max}[p] = \max\!\big(\texttt{running\_max}[p], \, |\nabla_p \mathcal{L}|\big).
    \end{equation*}
  \end{enumerate}
  \item Return \texttt{running\_max}.
\end{enumerate}

\subsection*{Code Illustration}
\begin{lstlisting}[style=pycode]
param_dict  = {n: p for n, p in model.named_parameters() if p.requires_grad}
running_max = {n: torch.zeros_like(p, dtype=torch.float32) for n, p in param_dict.items()}

for step, (I, Y) in enumerate(batches, 1):
    model.zero_grad(set_to_none=True)
    loss = model(I, labels=I).loss  # Hugging Face auto-aligns
    loss.backward()
    for name, p in param_dict.items():
        if p.grad is not None:
            running_max[name] = torch.maximum(running_max[name], p.grad.detach().abs())
    if step >= MAX_STEPS:
        break
\end{lstlisting}

\section{Deepseek-R1-Distill-Qwen-1.5B Gradient Scan Results}
The top-10 weights in the model (excluding embedding weights and layer-norm weights):
\begin{table}[H]
\centering
\begin{tabular}{|c|c|c|c|c|}
\hline
Local Rank & Global Rank & Weight & Element & $|\nabla \mathcal{L}|$ \\
\hline
1  & 2  & model.layers.3.mlp.gate\_proj.weight & (8521, 713) & $1.480\times 10^{0}$ \\
2  & 6  & lm\_head.weight                      & (1154, 507) & $1.152\times 10^{0}$ \\
3  & 9  & model.layers.3.mlp.gate\_proj.weight & (8521, 872) & $9.546\times 10^{-1}$ \\
4  & 10 & lm\_head.weight                      & (11,   507) & $9.478\times 10^{-1}$ \\
5  & 11 & model.layers.1.mlp.gate\_proj.weight & (6988, 147) & $9.453\times 10^{-1}$ \\
6  & 12 & model.layers.0.self\_attn.v\_proj.weight & (98, 900) & $9.424\times 10^{-1}$ \\
7  & 13 & lm\_head.weight                      & (1154, 1283) & $9.375\times 10^{-1}$ \\
8  & 17 & model.layers.3.mlp.gate\_proj.weight & (8521, 520) & $9.028\times 10^{-1}$ \\
9  & 19 & lm\_head.weight                      & (11, 1283) & $8.203\times 10^{-1}$ \\
10 & 21 & model.layers.0.mlp.down\_proj.weight & (147, 2385) & $8.120\times 10^{-1}$ \\
\hline
\end{tabular}
\end{table}

\end{document}
